\chapter{XML Basics}

\section{What is XML?}
XML is the \emph{Extensible Markup Language}. It is designed to
improve the functionality of the Web by providing more flexible and
adaptable information identification.

It is called extensible because it is not a fixed format like HTML (a single,
predefined markup language). Instead, XML is actually a \emph{metalanguage}
-- a language for describing other languages -- which lets you
design your own customized markup languages for limitless different
types of documents. XML can do this because it's written in \emph{SGML}, the
international standard metalanguage for text markup systems (ISO 8879).



\section{XML: Data Elements}
In XML a single data item is enclosed by two tags:
\begin{verbatim}
  <node_name>ground</node_name>
\end{verbatim}
an \emph{opening tag} and a \emph{closing tag}. The whole assembly is called an element.
In XML, an element can have

\begin{center}  
  \begin{tabular}{|l|l|}
    \hline
    a value    & \verb!<node>ground</node>!\\
    \hline 
    attributes & \verb!<node name="ground"></node>!\\
    \hline 
    children &  
    \begin{minipage}{10cm}
\begin{verbatim}
<nodelist>
  <node1>....</node1>
  <node2>....</node2>
</nodelist>
\end{verbatim}
    \end{minipage}\\
    \hline
  \end{tabular}
\end{center}


or a combination of the above.



\section{XML: Empty Elements}


If an element is empty, i.e.~it contains neither a value nor children,

\verb!<problem name="demo3d"></problem>!

\noindent
the opening and the closing tag can be combined into one:

\verb!<problem name="demo3d" />!